\documentclass[11pt]{article}
\usepackage[margin=2.2cm]{geometry}
\usepackage{amsmath,amssymb,mathtools}
\usepackage[shortlabels]{enumitem}
\setlist[enumerate,1]{label=\textbf{\arabic*.}, leftmargin=*, itemsep=10pt, topsep=6pt}
\usepackage{fancyhdr}
\usepackage{xcolor}

\newcommand{\R}{\mathbb{R}}
\newcommand{\E}{\mathbb{E}}
\DeclareMathOperator{\Var}{Var}
\DeclareMathOperator{\Cov}{Cov}
\newcommand{\ans}[1]{\par\noindent\textbf{\color{blue!50!black}Answer: #1}}

\pagestyle{fancy}
\fancyhf{}
\fancyhead[L]{Mathematical Finance \& Stochastic Processes}
\fancyhead[R]{Practice Quiz C --- Solutions}
\fancyfoot[C]{\thepage}
\renewcommand{\headrulewidth}{0.4pt}

\begin{document}

\begin{center}
{\LARGE\bfseries Practice Quiz C --- Solutions}
\end{center}
\vspace{6pt}
\hrule
\vspace{10pt}

\begin{enumerate}

\item[\textbf{Q1.}] \ans{A, B, D.} By the Toolbox rule for building new stopping times (Result 3): minima, maxima, and sums of stopping times are stopping times, and the minimum of a stopping time and a deterministic constant is a stopping time.
\begin{itemize}[nosep]
\item[A] $T_1\wedge T_2$: minimum of two stopping times. A stopping time.
\item[B] $T_1\vee T_2$: maximum of two stopping times. A stopping time.
\item[C] $T_2-T_1$: a \emph{difference} of stopping times is \textbf{not} covered by the rule, and is not a stopping time in general -- deciding $\{T_2-T_1\le t\}$ can require comparing two not-yet-fully-revealed random times, which is not $\mathcal F_t$-measurable information.
\item[D] $(T_1\wedge4)+(T_2\wedge4)$: each summand is the minimum of a stopping time and a constant, hence a stopping time; their sum is again a stopping time.
\end{itemize}

\item[\textbf{Q2.}] \ans{A.} By It\^o's formula, with $f(x,t)=x^2+at$, $f_x=2x$, $f_{xx}=2$, $f_t=a$:
\[
dY_t=f_x(B_t,t)\,dB_t+\Big(f_t(B_t,t)+\tfrac12f_{xx}(B_t,t)\Big)dt=2B_t\,dB_t+(a+1)\,dt.
\]
$Y$ is a martingale iff the drift vanishes, i.e.\ $a+1=0$, so $a=-1$. (This recovers the familiar fact that $B_t^2-t$ is a martingale.)

\item[\textbf{Q3.}] \ans{A.} Since the components $B_t^{(1)},B_t^{(2)}$ are independent,
\[
\Cov(\lambda\cdot\boldsymbol B_t,\mu\cdot\boldsymbol B_s)=\sum_{i,j}\lambda_i\mu_j\Cov(B_t^{(i)},B_s^{(j)})=\sum_i\lambda_i\mu_i\Cov(B_t^{(i)},B_s^{(i)})=(\lambda\cdot\mu)(s\wedge t),
\]
using that $\Cov(B_t^{(i)},B_s^{(j)})=0$ for $i\ne j$ (independent coordinates) and $=s\wedge t$ for $i=j$. With $s\le t$ and $\lambda\cdot\mu=\rho$, this is $\rho s$.

\item[\textbf{Q4.}] \ans{A.} Since $M_t=W_tB_t-\rho t$ is a martingale, the product rule $d(W_tB_t)=W_t\,dB_t+B_t\,dW_t+d\langle W,B\rangle_t$ forces the finite-variation part of $d(W_tB_t)$ to cancel against $\rho\,dt$, i.e.\ $\langle W,B\rangle_t=\rho t$. Hence $dM_t=W_t\,dB_t+B_t\,dW_t$, and by bracket bilinearity,
\[
\langle M,M\rangle_t=\int_0^tW_s^2\,ds+\int_0^tB_s^2\,ds+2\rho\int_0^tW_sB_s\,ds.
\]
Taking expectations, using $\E[W_s^2]=\E[B_s^2]=s$ and $\E[W_sB_s]=\Cov(W_s,B_s)=\rho s$:
\[
\E[\langle M,M\rangle_t]=\int_0^t\big(s+s+2\rho\cdot\rho s\big)\,ds=2(1+\rho^2)\int_0^ts\,ds=(1+\rho^2)t^2.
\]
With $\rho=0.5$: $(1+0.25)t^2=1.25\,t^2$.

\item[\textbf{Q5.}] \ans{A.} By Result 17 (Toolbox), $\Var(I_t)=t^3/3$, so $I_3\sim N(0,9)$, i.e.\ $I_3=3Z$ for $Z\sim N(0,1)$. Hence $P(I_3>3)=P(Z>1)=1-\Phi(1)$.

\item[\textbf{Q6.}] \ans{A.} With $f(x)=x^4$, $f'(x)=4x^3$, $f''(x)=12x^2$, It\^o's formula gives
\[
d(B_t^4)=4B_t^3\,dB_t+\tfrac12\cdot12B_t^2\,dt=4B_t^3\,dB_t+6B_t^2\,dt.
\]
Hence $dX_t=d(B_t^4)-6B_t^2\,dt=4B_t^3\,dB_t$: the drift term cancels exactly, leaving a pure stochastic integral. So $X$ is a martingale. (This generalises the familiar fact that $B_t^2-t$ is a martingale.)

\item[\textbf{Q7.}] \ans{A.} By the It\^o isometry, $\E[M_3^2]=\int_0^3s^2\,ds=9$. By Doob's $L^2$ maximal inequality (Result 19), $\E\big[\sup_{t\in[0,3]}M_t^2\big]\le4\,\E[M_3^2]=4\cdot9=36$.

\item[\textbf{Q8.}] \ans{A.} Expand $(x+2)^4=x^4+8x^3+24x^2+32x+16$ and substitute $x=B_t$, using $\E[B_t]=0$, $\E[B_t^2]=t$, $\E[B_t^3]=0$, $\E[B_t^4]=3t^2$:
\[
\E[(B_t+2)^4]=3t^2+8\cdot0+24t+32\cdot0+16=3t^2+24t+16.
\]

\item[\textbf{Q9.}] \ans{A.} By bracket bilinearity and the orthogonality of $B^{(1)}$ and $B^{(2)}$ (Problem 5.4), the cross term vanishes:
\[
\langle X,X\rangle_t=\int_0^t\cos^2(s)\,ds+\int_0^t\sin^2(s)\,ds=\int_0^t\big(\cos^2(s)+\sin^2(s)\big)\,ds=\int_0^t1\,ds=t.
\]
Since $\langle X,X\rangle_t=t$ is deterministic, continuous and increasing, L\'evy's characterisation identifies $X$ as a standard Brownian motion.

\item[\textbf{Q10.}] \ans{A.} $W_2=2B_{1/2}$, so $\Cov(B_1,W_2)=2\Cov(B_1,B_{1/2})=2\min(1,0.5)=2\cdot0.5=1$.

\item[\textbf{Q11.}] \ans{A, C.}
\begin{itemize}[nosep]
\item[A] $B_s^3$ is $\mathcal F_s$-measurable and $B_t-B_s$ is independent of $\mathcal F_s$, so $\E[B_s^3(B_t-B_s)]=\E[B_s^3]\cdot\E[B_t-B_s]=0\cdot0=0$.
\item[B] Write $B_t=B_s+(B_t-B_s)$ and expand $B_t^3(B_t-B_s)=B_t^4-B_t^3B_s$. Now $\E[B_t^4]=3t^2$, and by Wick's theorem (for jointly Gaussian mean-zero $(B_t,B_s)$), $\E[B_t^3B_s]=3\Var(B_t)\Cov(B_t,B_s)=3ts$. Hence $\E[B_t^3(B_t-B_s)]=3t^2-3ts=3t(t-s)\ne0$ in general. \textbf{Not} zero -- despite superficially resembling A, the factor $B_t^3$ is \emph{not} independent of $B_t-B_s$ (both involve $B_t$), so the naive ``independent factor times mean-zero increment'' shortcut does not apply here.
\item[C] $(B_t-B_s)^3$ is (up to sign) the cube of a centered Gaussian random variable $B_t-B_s\sim N(0,t-s)$: all odd central moments vanish, so $\E[(B_t-B_s)^3]=0$.
\item[D] $B_s^2$ is $\mathcal F_s$-measurable and $(B_t-B_s)^2$ is a function of the increment, independent of $\mathcal F_s$: $\E[B_s^2(B_t-B_s)^2]=\E[B_s^2]\cdot\E[(B_t-B_s)^2]=s(t-s)\ne0$ in general. \textbf{Not} zero.
\end{itemize}

\item[\textbf{Q12.}] \ans{A.} This is the general conditional-MGF identity for a continuous local martingale with deterministic quadratic variation $g$ (Problem 6.2(b)(ii)):
\[
\E\big[\exp(\lambda(M_t-M_s))\mid\mathcal F_s\big]=\exp\Big(\tfrac{\lambda^2}2\big(g(t)-g(s)\big)\Big).
\]
With $g(t)=t^2$: $g(t)-g(s)=t^2-s^2$, giving $\exp\big(\tfrac{\lambda^2}2(t^2-s^2)\big)$.

\item[\textbf{Q13.}] \ans{A.} $C_t=Z_t^2=\exp\big(2\lambda\cdot\boldsymbol B_t-|\lambda|^2t\big)$. Since the coordinates of $\boldsymbol B_t$ are independent $N(0,t)$, for any $v\in\R^d$, $\E[\exp(v\cdot\boldsymbol B_t)]=\exp\big(\tfrac12|v|^2t\big)$. With $v=2\lambda$: $\E[\exp(2\lambda\cdot\boldsymbol B_t)]=\exp(2|\lambda|^2t)$. Hence
\[
\E[C_t]=\exp(2|\lambda|^2t)\cdot e^{-|\lambda|^2t}=e^{|\lambda|^2t}.
\]
(Consistent with $C=Z^2$ being a submartingale, since $x\mapsto x^2$ is convex and $Z$ is a martingale: $\E[C_t]=e^{|\lambda|^2t}\ge1=\E[C_0]$ and is strictly increasing.)

\item[\textbf{Q14.}] \ans{A.} By the It\^o isometry, $\E[X_s^2]=\int_0^se^{8u}\,du=\tfrac18(e^{8s}-1)$. Applying the isometry again to the outer integral,
\[
\E\left[\left(\int_0^tX_s\,dB_s\right)^2\right]=\int_0^t\E[X_s^2]\,ds=\int_0^t\tfrac18(e^{8s}-1)\,ds=\tfrac1{64}(e^{8t}-1)-\tfrac t8.
\]

\item[\textbf{Q15.}] \ans{D.} $X_t=e^{-4t}B_{e^{8t}}$ is a deterministic rescaling of Brownian motion at a deterministic time, hence centered Gaussian: $\E[X_t]=e^{-4t}\E[B_{e^{8t}}]=0$. For $0\le s<t$ (so $e^{8s}<e^{8t}$),
\[
\Cov(X_s,X_t)=e^{-4s-4t}\Cov(B_{e^{8s}},B_{e^{8t}})=e^{-4s-4t}\min(e^{8s},e^{8t})=e^{-4s-4t}e^{8s}=e^{4s-4t}=e^{-4(t-s)}.
\]

\item[\textbf{Q16.}] \ans{B.} $e^{\lambda X_t}=e^{\lambda B_t+4\lambda t}=e^{\lambda B_t-\lambda^2t/2}\cdot e^{(4\lambda+\lambda^2/2)t}$. Since $e^{\lambda B_t-\lambda^2t/2}$ is always a martingale, the product is a martingale iff the remaining deterministic exponential factor is constant, i.e.\ $4\lambda+\lambda^2/2=0$, giving $\lambda=0$ or $\lambda=-8$. So $e^{-8X_t}$ is a martingale.

\end{enumerate}

\end{document}
